\chap{Введение}

Цели данной лабораторной работы:
\begin{enumerate}
	\item получить представление о составе дистрибутивов Linux;
	\item приобрести практические навыки установки и запуска дистрибутива Linux в виртуальной машине Oracle VM VirtualBox, получения справочной информации о системе и установленных приложениях с помощью интерфейса командной строки.
\end{enumerate}

\chap{Выполнение}

\sect{Задание 1}

\begin{figure}[!htbp]
	\centering
	\includegraphics[width=0.6\linewidth]{"Screenshot From 2026-09-09 13-22-47"}
	\caption{Терминал запущенной виртуальной машины}
\end{figure}
\FloatBarrier

\sect{Задание 2}

Обновление списка пакетов. 

\begin{figure}[!htbp]
	\centering
	\includegraphics[width=0.6\linewidth]{"Screenshot From 2026-09-09 13-22-01"}
	\caption{Список пакетов, имеющих обновление (update)}
\end{figure}
\FloatBarrier

%Занесите в отчет количество пакетов, подлежащих обновлению, сколько места для этого потребуется? Определите, какие репозитории используются для обновления. Объясните, почему, прежде чем выполнять фактическое обновление, важно проанализировать информацию о репозитории и размере обновления.

Для обновления доступно 9 пакетов, требуется 3Мб места на диске. Прежде чем выполнять фактическое обновление, важно проанализировать информацию о репозитории и размере обновления, так как для скачивания и установки используются ресурсы системы, а без них обновление невозможно.

\begin{figure}[!htbp]
	\centering
	\includegraphics[width=0.6\linewidth]{"Screenshot From 2026-09-09 13-22-32"}
	\caption{Список пакетов, имеющих обновление (upgrade)}
\end{figure}
\FloatBarrier

Обновление системы.
\begin{figure}[!htbp]
	\centering
	\includegraphics[width=0.6\linewidth]{"Screenshot From 2026-09-09 13-24-24"}
	\caption{Подтверждение обновления системы}
\end{figure}
\FloatBarrier

Вывод информации о дистрибутиве: версия операционной системы, версия ядра. 
\begin{figure}[!htbp]
	\centering
	\includegraphics[width=0.6\linewidth]{"Screenshot From 2026-09-21 15-35-04"}
	\caption{Информация о дистрибутиве}
\end{figure}
\FloatBarrier

Вывод информации о модели процессора, количестве ядер процессора, объеме оперативной памяти.
\clearpage
\begin{figure}[!htbp]
	\centering
	\includegraphics[width=0.6\linewidth]{"Screenshot From 2026-09-09 13-32-40"}
	\includegraphics[width=0.6\linewidth]{"Screenshot From 2026-09-09 13-33-13"}
	\caption{Информация о технических характеристиках устройства}
\end{figure}
\FloatBarrier

Вывод команды uname -a. 
\begin{figure}[!htbp]
	\centering
	\includegraphics[width=0.6\linewidth]{"Screenshot From 2026-09-21 15-23-42"}
	\caption{Вывод команды uname -a}
\end{figure}
\FloatBarrier
%В отчете приведите разбор полученного вывода.
Разбор полученного вывода:
\begin{itemize}
	\item Linux -- тип ядра;
	\item elanskaya -- имя пользователя;
	\item 7.0.0 -- версия ядра;
	\item 31 -- номер сборки;
	\item generic -- тип сборки;
	\item \#31-Ubuntu -- номер сборки;
	\item SMP -- поддержка многопроцессорности;
	\item PREEMPT\_DYNAMIC -- режим вытеснения;
	\item Sat Aug  1 04:26:38 UTC 2026 -- дата компиляции ядра;
	\item x86\_64 -- тип архитектуры;
	\item GNU/Linux -- тип ОС.
\end{itemize}

Вывод системы инициализации используется в этом дистрибутиве Linux.
\begin{figure}[!htbp]
	\centering
	\includegraphics[width=0.6\linewidth]{"Screenshot From 2026-09-09 13-48-39"}
	\caption{Команда вывода названия системы инициализации}
\end{figure}
\FloatBarrier
Список доступных репозиториев для системы.
\begin{figure}[!htbp]
	\centering
	\includegraphics[width=0.6\linewidth]{"Screenshot From 2026-09-23 12-57-10"}
	\caption{Список доступных репозиториев}
\end{figure}
\FloatBarrier

Установка на виртуальную машину файлового менеджера Midnight Commander и текстового редактора Vim.
\begin{figure}[!htbp]
	\centering
	\includegraphics[width=0.6\linewidth]{"Screenshot From 2026-09-09 13-38-45"}
	\includegraphics[width=0.6\linewidth]{"Screenshot From 2026-09-09 13-40-19"}
	\caption{Установка Vim и MC}
\end{figure}
\FloatBarrier

Список установленных программ (Midnight Commander и Vim в нем появились).
\begin{figure}[!htbp]
	\centering
	\includegraphics[width=0.6\linewidth]{"Screenshot From 2026-09-21 16-00-11"}
	\caption{Список установленных программ}
\end{figure}
\FloatBarrier

\begin{figure}[!htbp]
	\centering
	\includegraphics[width=0.6\linewidth]{"Screenshot From 2026-09-21 15-18-02"}
	\caption{Открытый текстовый редактор Vim}
\end{figure}
\FloatBarrier

Получение информации, из какого репозитория был установлен пакет.
\begin{figure}[!htbp]
	\centering
	\includegraphics[width=0.6\linewidth]{"Screenshot From 2026-09-23 13-06-46"}
	\caption{Вывод названия репозитория}
\end{figure}
\FloatBarrier

Место расположения исполняемых файлов установленного программного обеспечения.
\begin{figure}[!htbp]
	\centering
	\includegraphics[width=0.6\linewidth]{"Screenshot From 2026-09-23 13-00-11"}
	\caption{Место расположения исполняемых файлов}
\end{figure}
\FloatBarrier

Вывод зависимостей Midnight Commander.
\begin{figure}[!htbp]
	\centering
	\includegraphics[width=0.6\linewidth]{"Screenshot From 2026-09-09 13-50-34"}
	\caption{Вывод зависимостей MC}
\end{figure}
\FloatBarrier

\sect{Дополнительное задание}

В качестве DEB-дистрибутива расссматривалась Ubuntu 26.04, в качестве RPM-дистрибутива -- Fedora 44.

\begin{figure}[!htbp]
	\centering
	\includegraphics[width=0.6\linewidth]{"Screenshot From 2026-09-23 14-00-27"}
	\caption{Рабочий стол Fedora}
\end{figure}
\FloatBarrier

\begin{table}[!htbp]
	\caption{Сравнение дистрибутивов}
	\begin{tabularx}{\textwidth}{|p{7cm}|X|X|}
		\hline
		& \textbf{DEB-дистрибутив} & \textbf{RPM-дистрибутив} \\\hline
		Пакетный менеджер (основной) & apt & dnf/yum \\\hline
		Низкоуровневый пакетный менеджер & dpkg & rpm\\\hline
		Формат пакетов & .deb & .rpm \\\hline
		Обновление системы & sudo apt update & sudo dnf upgrade\\\hline
		Установка пакета & sudo apt install <pkg> & sudo dnf install <pkg>\\\hline
		Список установленных пакетов & dpkg -l & rpm -qa\\\hline
		Список репозиториев & /etc/apt/sources.list & /etc/yum.repos.d\\\hline
		Зависимости пакета & apt show <pkg> & dnf info <pkg>\\\hline
	\end{tabularx}
\end{table}

%\lstinputlisting[style=python, caption={}]{}

\chap{Вывод}

В ходе лабораторной работы:
\begin{itemize}
	\item было получено представление о составе дистрибутивов Linux
	\item были приобретены практические навыки установки и запуска; дистрибутива Linux в виртуальной машине Oracle VM VirtualBox, получения справочной информации о системе и установленных приложениях с помощью интерфейса командной строки.
\end{itemize} 

\chap{Контрольные вопросы}

\begin{enumerate}
	\item \textit{Что такое дистрибутив? Состав дистрибутива.}\newline 
	Форма распространения системного программного обеспечения. Обычно содержит программу для начальной инициализации системы, программу-установщик и набор специальных файлов, содержащих отдельные части системы.
	\item \textit{Функции ядра операционной системы.}\newline 
	Взаимодействие с аппаратным обеспечением, распределение оперативной памяти, управление процессами, работа с файловыми системами.
	\item \textit{Классическая и микроядерная архитектуры ОС.}\newline 
	У классической больший функционал, высокая производительность, но меньшая надежность; у микроядерной минимальный функционал, меньшая производительность, но более высокая надежность благоларя изолированности.
	\item \textit{Дайте определение, что такое репозиторий. Что в них хранится?}\newline 
	Централизованные, доверенные источники ПО, обеспечивающие безопасность и стабильность. Хранит пакеты и их индексы.
	\item \textit{Как получить список доступных сторонних репозиториев?}\newline 
	Просмотреть содержимое sources.list.
	\item \textit{Назначение и функции пакетных менеджеров apt и dpkg.}\newline 
	Apt пакетный менеджер верхнего уровня, имеет выход в сеть для проверки и скачивания обновлений; dpkg пакетный менеджер низкого уровня, устанавливает пакеты.
	\item \textit{Что происходит при выполнении apt update и apt upgrade?}\newline 
	Apt update проверяет наличие обновлений в удаленном репозитории, apt upgrade скачивает и уставливает их.
	\item \textit{Что такое зависимости пакета и как APT их разрешает?}\newline 
	Зависимости (depends) -- это другие пакеты, необходимые для работы данного пакета. Apt читает метаданные пакета, проверяет установленные версии и загружает недостающие необходимые пакеты.
\end{enumerate}
